% Topic 1.2.04 — Independence of Events.

\section{Независимость событий}
\label{sec:1.2.04-independence}

\begin{ticket}
Независимость событий. Попарная независимость и независимость в совокупности.
\end{ticket}

\subsection*{Теория}

Независимость~--- предположение о том, что одна порция информации не несёт в себе никаких вероятностных сведений о другой; формально это означает, что совместная вероятность факторизуется. Определение чисто мультипликативное: ни механизм взаимодействия, ни причинно-следственная интерпретация в нём не заложены.

\begin{definition}[Два события]
\label{def:indep-two}
События $A, B \in \mathcal{F}$ называются \emph{независимыми}, если
\[
  \Prob(A \cap B) \;=\; \Prob(A)\, \Prob(B).
\]
\end{definition}

\begin{proposition}[Условная переформулировка]
\label{prop:indep-conditional}
Если $\Prob(B) > 0$, то~$A$ и~$B$ независимы тогда и только тогда, когда $\Prob(A \mid B) = \Prob(A)$. Симметрично, при $\Prob(A) > 0$~--- тогда и только тогда, когда $\Prob(B \mid A) = \Prob(B)$.
\end{proposition}
\proofof{indep-conditional}

\begin{proposition}[Сохранение независимости при переходе к дополнению]
\label{prop:indep-complements}
Если события~$A, B$ независимы, то независимы и пары $(A, B^c)$, $(A^c, B)$, $(A^c, B^c)$.
\end{proposition}
\proofof{indep-complements}

Перенос понятия независимости с пары на семейство событий допускает две естественные интерпретации, которые \emph{не} эквивалентны.

\begin{definition}[Попарная независимость и независимость в совокупности]
\label{def:indep-many}
Пусть $\{A_i\}_{i \in I}$~--- (конечное или счётное) семейство событий.
\begin{itemize}
  \item Семейство называется \emph{попарно независимым}, если $\Prob(A_i \cap A_j) = \Prob(A_i) \Prob(A_j)$ для каждой пары $i \neq j$ из~$I$.
  \item Семейство называется \emph{независимым в совокупности} (\emph{независимыми в совокупности} или \emph{взаимно независимыми}), если для всякого конечного $J \subseteq I$
        \[
          \Prob\Bigl(\bigcap_{i \in J} A_i\Bigr) \;=\; \prod_{i \in J} \Prob(A_i).
        \]
\end{itemize}
Независимость в совокупности влечёт попарную независимость; обратное при $|I| \geq 3$ неверно, что демонстрирует пример Бернштейна (\cref{ex:bernstein}).
\end{definition}

Для независимости в совокупности факторизация требуется на \emph{всех} непустых конечных подмножествах, а не только на полном пересечении. Именно это свойство делает независимость в совокупности достаточно сильной, чтобы переноситься на операции над событиями.

\begin{proposition}[Устойчивость независимости в совокупности]
\label{prop:indep-stability}
Пусть события $A_1, \dots, A_n$ независимы в совокупности. Тогда:
\begin{enumerate}[label=(\alph*)]
  \item замена любого~$A_i$ на $A_i^c$ сохраняет независимость в совокупности;
  \item при $1 \leq r < n$ события $B = A_1 \cap \cdots \cap A_r$ и $C = A_{r+1} \cap \cdots \cap A_n$ независимы (``блочная'' независимость).
\end{enumerate}
\end{proposition}
\proofof{indep-stability}

\medskip
\noindent\textbf{Независимость случайных величин.} Понятие независимости переносится на случайные величины через прообразы борелевских множеств.

\begin{definition}
\label{def:indep-rv}
Случайные величины $X_1, \dots, X_n$ на $(\Omega, \mathcal{F}, \Prob)$ называются \emph{независимыми (в совокупности)}, если для любых борелевских множеств $B_1, \dots, B_n \in \mathcal{B}(\R)$
\[
  \Prob(X_1 \in B_1, \dots, X_n \in B_n) \;=\; \prod_{i=1}^{n} \Prob(X_i \in B_i).
\]
Эквивалентным образом~--- в терминах совместной и одномерных функций распределения~--- $F_{(X_1, \dots, X_n)}(x_1, \dots, x_n) = \prod_i F_{X_i}(x_i)$ для всех $x_i \in \R$.
\end{definition}

\begin{theorem}[Мультипликативность математического ожидания]
\label{thm:indep-multiplicative}
Если $X, Y \in L^1$ независимы, то $XY \in L^1$ и
\[
  \E(XY) \;=\; \E X \cdot \E Y.
\]
Как следствие, независимые величины из $L^2$ некоррелированы: $\Cov(X, Y) = 0$. Обратное в общем случае неверно (\cref{ex:uncorrelated-dependent}).
\end{theorem}
\proofof{indep-multiplicative}

\begin{corollary}[Дисперсия суммы независимых величин]
\label{cor:var-sum-indep}
Если $X_1, \dots, X_n \in L^2$ независимы в совокупности, то
\[
  \Var\Bigl(\sum_{i=1}^{n} X_i\Bigr) \;=\; \sum_{i=1}^{n} \Var X_i.
\]
\end{corollary}
\proofof{var-sum-indep}

\begin{remark}
Независимость~--- рабочая предпосылка предельных теорем (\cref{sec:1.2.05-limit-theorems}), статистических моделей выборок (\cref{sec:1.2.10-chi2-f-student}), а также формализма цепочек условных вероятностей (\cref{sec:1.2.09-markov-chains}), где она ослабляется до марковского свойства. Полная теоретико-мерная конструкция независимости $\sigma$-алгебр и связанные с ней теоремы о продолжении изложены в~\cite[гл.~II, \S~5]{Shiryaev1989}.
\end{remark}

\subsection*{Доказательства}

\begin{proof}[\Cref{prop:indep-conditional}]
\proofanchor{indep-conditional}
По определению $\Prob(A \mid B) = \Prob(A \cap B) / \Prob(B)$. При $\Prob(B) > 0$ равенство правой части~$\Prob(A)$ эквивалентно $\Prob(A \cap B) = \Prob(A) \Prob(B)$~--- определению независимости. Роли~$A$ и~$B$ в факторизованной форме симметричны.
\end{proof}

\begin{proof}[\Cref{prop:indep-complements}]
\proofanchor{indep-complements}
Разложим~$A$ в дизъюнктное объединение $A = (A \cap B) \sqcup (A \cap B^c)$; тогда
\[
  \Prob(A \cap B^c) = \Prob(A) - \Prob(A \cap B) = \Prob(A) - \Prob(A) \Prob(B) = \Prob(A) (1 - \Prob(B)) = \Prob(A) \Prob(B^c),
\]
где использованы~\cref{def:indep-two} и правило дополнения~\cref{prop:prob-elementary}\,(c). Значит, события~$A$ и~$B^c$ независимы. Применяя это рассуждение дважды (сначала к~$A, B$, затем к~$B^c, A$), получим независимость остальных двух пар.
\end{proof}

\begin{proof}[\Cref{prop:indep-stability}]
\proofanchor{indep-stability}
\emph{(a)} Достаточно показать, что замена~$A_n$ на~$A_n^c$ сохраняет независимость в совокупности; повторное применение даёт общий случай. Для произвольного $J \subseteq \{1, \dots, n-1\}$ применим независимость в совокупности к множествам $J \cup \{n\}$ и~$J$:
\[
  \Prob\Bigl(\bigcap_{i \in J} A_i \cap A_n^c\Bigr) = \Prob\Bigl(\bigcap_{i \in J} A_i\Bigr) - \Prob\Bigl(\bigcap_{i \in J} A_i \cap A_n\Bigr) = \prod_{i \in J} \Prob(A_i)\, (1 - \Prob(A_n)),
\]
а это в точности факторизация для семейства, в котором~$A_n$ заменено на~$A_n^c$.

\emph{(b)} По независимости в совокупности $A_1, \dots, A_n$, применённой к полному множеству индексов $\{1, \dots, n\}$,
\[
  \Prob(B \cap C) = \Prob\Bigl(\bigcap_{i=1}^{n} A_i\Bigr) = \prod_{i=1}^{n} \Prob(A_i) = \prod_{i=1}^{r} \Prob(A_i) \cdot \prod_{i=r+1}^{n} \Prob(A_i) = \Prob(B)\, \Prob(C),
\]
где в третьем равенстве используется независимость в совокупности на префиксном и суффиксном подмножествах для вычисления~$\Prob(B)$ и~$\Prob(C)$ по отдельности.
\end{proof}

\begin{proof}[\Cref{thm:indep-multiplicative}]
\proofanchor{indep-multiplicative}
Сначала рассмотрим $X = \mathbf{1}_A$, $Y = \mathbf{1}_B$ для событий $A, B$. По независимости порождённых~$X$ и~$Y$ $\sigma$-алгебр (что в случае индикаторов сводится к~\cref{def:indep-two}) $\E(XY) = \E \mathbf{1}_{A \cap B} = \Prob(A \cap B) = \Prob(A) \Prob(B) = \E X \cdot \E Y$. По линейности (\cref{prop:exp-properties}\,(a)) тождество распространяется на простые функции $X = \sum_i a_i \mathbf{1}_{A_i}$, $Y = \sum_j b_j \mathbf{1}_{B_j}$ с $\{A_i\}$ и $\{B_j\}$~--- разбиениями, ``развязанными'' независимостью соответствующих $\sigma$-алгебр. Общий случай $X, Y \in L^1$ получается аппроксимацией неотрицательных $X, Y \geq 0$ простыми неотрицательными функциями $X_n \uparrow X, Y_n \uparrow Y$ и применением теоремы о монотонной сходимости (это законно, поскольку $X_n Y_n \leq XY$ поточечно и $\E X_n Y_n = \E X_n \cdot \E Y_n \to \E X \cdot \E Y$). Разложение $X = X^+ - X^-$, $Y = Y^+ - Y^-$ переносит результат на знакопеременный случай. Полное изложение~--- в~\cite[гл.~II, \S~6]{Shiryaev1989}.

Следствие о ковариации немедленно: $\Cov(X, Y) = \E XY - \E X\, \E Y = 0$.
\end{proof}

\begin{proof}[\Cref{cor:var-sum-indep}]
\proofanchor{var-sum-indep}
По~\cref{prop:cov-properties}\,(e) $\Var(\sum X_i) = \sum \Var X_i + 2 \sum_{i < j} \Cov(X_i, X_j)$. Пары $(X_i, X_j)$ при $i \neq j$ попарно независимы (как проекция независимости в совокупности на два индекса), а потому некоррелированы по~\cref{thm:indep-multiplicative}, и недиагональная сумма обнуляется.
\end{proof}

\subsection*{Примеры}

\begin{example}[Два бросания монеты]
\label{ex:two-coins}
Бросаем правильную монету дважды независимо; пусть $A = \{\text{первый бросок~--- ``орёл''}\}$ и $B = \{\text{второй бросок~--- ``орёл''}\}$. Возьмём $\Omega = \{\text{О}, \text{Р}\}^2$ с классической мерой. Тогда $\Prob(A) = \Prob(B) = 1/2$ и $\Prob(A \cap B) = 1/4 = \Prob(A) \Prob(B)$, так что~$A$ и~$B$ независимы~--- проверка того, что построенная модель согласуется с интуитивным понятием. Эквивалентно, $\Prob(A \mid B) = 1/2 = \Prob(A)$ (\cref{prop:indep-conditional}).
\end{example}

\begin{example}[Бернштейн: попарная, но не совокупная независимость]
\label{ex:bernstein}
Возьмём $\Omega = \{1, 2, 3, 4\}$ с равномерной мерой и события $A = \{1, 2\}$, $B = \{1, 3\}$, $C = \{1, 4\}$. Каждое из них имеет вероятность~$1/2$, всякое попарное пересечение равно $\{1\}$ с вероятностью $1/4 = (1/2)(1/2)$, так что семейство попарно независимо. Однако $A \cap B \cap C = \{1\}$, и $\Prob(A \cap B \cap C) = 1/4 \neq 1/8 = (1/2)^3$. Тройка \emph{не} независима в совокупности~--- классическая иллюстрация того, что дополнительные равенства из~\cref{def:indep-many} нельзя свести к одной только попарной проверке.
\end{example}

\begin{example}[Число выпадений ``орла'', продолжение]
\label{ex:binomial-indep}
В обозначениях $S_n = X_1 + \cdots + X_n$ из~\cref{ex:binomial-moments} величины~$X_i$ независимы в совокупности (по построению независимых испытаний); тогда из~\cref{cor:var-sum-indep} следует $\Var S_n = \sum_i \Var X_i = n p(1 - p)$. Прежний вывод в~\cref{ex:binomial-moments} опирался лишь на попарную некоррелированность~--- строго более слабое предположение, которого, впрочем, оказалось достаточно для формулы дисперсии суммы.
\end{example}

\begin{example}[Из некоррелированности не следует независимость]
\label{ex:uncorrelated-not-indep}
Пара $X \sim \mathrm{Uniform}\{-1, 0, 1\}$ и $Y = X^2$ из~\cref{ex:uncorrelated-dependent} некоррелирована, но связана детерминированно, так что независимой не является: например, $\Prob(X = 0, Y = 0) = \Prob(X = 0) = 1/3$, тогда как $\Prob(X = 0)\, \Prob(Y = 0) = (1/3)(1/3) = 1/9$.
\end{example}

\begin{problem}
\label{prob:exactly-k-out-of-n}
Правильную игральную кость бросают~$n$ раз независимо. Найдите вероятность того, что число выпадений шестёрки равно~$k$, и проверьте, что полученное распределение в сумме по $k = 0, 1, \dots, n$ даёт~$1$.
\end{problem}

\begin{proof}[Решение]
Пусть $X_i = \mathbf{1}_{\{\text{бросок }i\text{ дал шестёрку}\}}$. Тогда $X_1, \dots, X_n$ независимы в совокупности, каждая с параметром $p = 1/6$. Сумма $S_n = \sum X_i$ считает шестёрки; событие $\{S_n = k\}$~--- объединение $\binom{n}{k}$ конфигураций с указанными $k$~``удачными'' бросками, причём на каждой такой конфигурации совместная вероятность факторизуется (независимость в совокупности) в $p^k (1-p)^{n-k}$. Дизъюнктное объединение даёт
\[
  \Prob(S_n = k) \;=\; \binom{n}{k} p^k (1 - p)^{n-k}, \qquad p = 1/6.
\]
По биномиальной теореме $\sum_{k=0}^n \binom{n}{k} p^k (1-p)^{n-k} = (p + (1 - p))^n = 1$. Это биномиальное распределение; оно вновь возникнет в~\cref{sec:1.2.06-standard-distributions} как одно из стандартных дискретных семейств.
\end{proof}
